\color{unidarkgreen}{\subsection{Логика деления стороны ВН (ЛД ВН)}}
\color{black}

\subsubsection{Общие сведения}

\begin{enumerate}[label=\arabic{section}.\arabic{subsection}.\arabic{subsubsection}.\arabic*,  align=left, labelsep=4pt, leftmargin=0pt, itemindent=57pt]

\begin{figure}[H]
\centering
\includegraphics[width=0.8\textwidth,height=0.8\textheight,keepaspectratio]{\fbpath/021.1791 ЛД ВН/img/LD_VN.pdf}
\captionof{figure}{Функциональная схема ЛД стороны ВН}\label{pic:LD_VN}
\end{figure}

\item
Ввод в работу логики деления высокой стороны трансформатора осуществляется с помощью
уставок \textbf{«ЛД ВН от ТЗНП»}. Оперативный ввод/вывод ЛД высокой стороны трансформатора
производится с помощью БУ (блока управления функциями) «Вывод ЛД ВН».   

\item
От сигнала пуска ТЗНП ВН, через выдержку времени \textbf{«Тср ТЗНП откл СВ(ШСВ) ВН»} логика
деления формирует сигналы на отключение СВ (ШСВ) ВН, а через выдержку времени \textbf{«Тср ТЗНП откл
В ВН»} логика деления формирует сигналы на отключение выключателей высокой стороны
трансформатора.


\small
\begin{longtable}{|>{\centering\arraybackslash}m{5.3cm}|>{\centering\arraybackslash}m{4.2cm}|>{\centering\arraybackslash}m{1cm}|>{\centering\arraybackslash}m{3.3cm}|>{\centering\arraybackslash}m{1.8cm}|}
% Первая страница: полный заголовок
\caption{Перечень уставок ЛД стороны ВН\hfill\vspace{-0.5\baselineskip}}\label{desc:tbl1}\\
\hline
\rowcolor{gray!30}
Наименование уставки & Диапазон & Шаг & Значение по умолчанию & Единица измерения \\
\hline
\endfirsthead

% Промежуточные страницы: «Продолжение таблицы»
\caption*{\hspace{3pt}\emph{Продолжение таблицы \ref{desc:tbl1}}\hfill\ }\vspace{-0.5\baselineskip} \\
\hline
\rowcolor{gray!30}
Наименование уставки & Диапазон & Шаг & Значение по умолчанию & Единица измерения \\
\hline
\endhead

%===m>ДТЗ|-
%===m
\end{longtable}
\normalsize	

\end{enumerate}